\section*{अध्याय \ref{ch:b2:randomvar} --- विविक्त यादृच्छिक चर}

\begin{solution}{exo:b2:randomvar:1}
\emph{द्विपद:} \omterm{def:b2:proba:independence}{स्वतंत्र} बर्नूली $X_i$ के साथ $X = \sum_{i=1}^n X_i$; $V(X_i) = \E(X_i^2) - \E(X_i)^2 = p - p^2$, और \omterm{def:b2:proba:independence}{स्वतंत्र}
चरों के \omterm{def:b2:randomvar:variance}{प्रसरण} जुड़ जाते हैं (\cref{thm:b2:randomvar:variancerules}):
\[
\E(X) = np, \qquad V(X) = np(1-p) .
\]
\emph{प्वासों:} $\E\bigl(X(X-1)\bigr) =
\sum_{k\geq2}k(k-1)e^{-\lambda}\frac{\lambda^k}{k!} = \lambda^2
e^{-\lambda}\sum_{j\geq0}\frac{\lambda^j}{j!} = \lambda^2$, अतः
\[
V(X) = \E(X^2) - \E(X)^2
= \lambda^2 + \lambda - \lambda^2 = \lambda .
\]
\emph{ज्यामितीय} ($q = 1 - p$): $\sum_{k\geq0}q^k = \frac{1}{1-q}$ का चक्रिका के भीतर दो बार अवकलन करने
पर (\cref{ch:b2:powerseries}) $\sum_{k\geq2}k(k-1)q^{k-2} =
\frac{2}{(1-q)^3}$, अतः
\[
\E\bigl(X(X-1)\bigr) = pq\sum_{k\geq2}k(k-1)q^{k-2}
= \frac{2q}{p^2},
\qquad
V(X) = \frac{2q}{p^2} + \frac1p - \frac{1}{p^2}
= \frac{q}{p^2} = \frac{1-p}{p^2} .
\]
\end{solution}

\begin{solution}{exo:b2:randomvar:2}
\emph{योग:} $n \in \N$ के लिए, असंयुक्तता तथा स्वतंत्रता से,
\[
\P(X + Y = n)
= \sum_{k=0}^n \P(X = k)\P(Y = n - k)
= e^{-(\lambda + \mu)}\frac{1}{n!}
\sum_{k=0}^n \binom nk \lambda^k\mu^{n-k}
= e^{-(\lambda+\mu)}\frac{(\lambda + \mu)^n}{n!}
\]
और द्विपद प्रमेय से: $X + Y \sim \mathcal{P}(\lambda + \mu)$।
\emph{सप्रतिबंध \omterm{def:b2:randomvar:law}{नियम}:} $0 \leq k \leq n$ के लिए,
\[
\P(X = k \mid X + Y = n)
= \frac{\P(X = k)\P(Y = n - k)}{\P(X + Y = n)}
= \binom nk
\Bigl(\frac{\lambda}{\lambda+\mu}\Bigr)^{k}
\Bigl(\frac{\mu}{\lambda+\mu}\Bigr)^{n-k} ,
\]
अर्थात् द्विपद \omterm{def:b2:randomvar:law}{नियम} $\mathcal{B}\bigl(n,
\frac{\lambda}{\lambda+\mu}\bigr)$: कुल गिनती दी हो तो हर \omterm{def:b2:proba:space}{घटना} \omterm{def:b2:proba:independence}{स्वतंत्र} रूप से
अपनी दर के समानुपाती प्रायिकता के साथ पहला स्रोत ``चुन'' लेती है।
\end{solution}

\begin{solution}{exo:b2:randomvar:3}
जब $k$ खिलौने अब भी शेष हों, तब हर नया डिब्बा प्रायिकता $\frac kn$ के साथ
कोई नया खिलौना लाता है, और यह अतीत से \omterm{def:b2:proba:independence}{स्वतंत्र} है: अगले नए खिलौने
के लिए प्रतीक्षा-समय $W_k$ ज्यामितीय $\mathcal{G}\bigl(\frac kn\bigr)$ है, जहाँ $\E(W_k) = \frac nk$, और $T_n = W_n + W_{n-1} + \dots + W_1$
(पहला डिब्बा सदा कोई नया खिलौना देता है: $W_n = 1$, जो $\E = n/n$ से संगत है)।
रैखिकता से,
\[
\E(T_n) = \sum_{k=1}^n \frac nk = n\sum_{k=1}^n\frac1k
\sim n\ln n ,
\]
जहाँ $\sum_{k\leq n}\frac1k = \ln n + \gamma + o(1)$ का उपयोग हुआ (\cref{ch:b2:comparison})। क़ीमत अंतिम कुछ खिलौनों को इकट्ठा
करने की है: आधे डिब्बे उसी अंतिम मुट्ठी भर पर जाते हैं।
\end{solution}

\begin{solution}{exo:b2:randomvar:4}
\omterm{def:b2:funcseq:def}{बिंदुवार} $X(\omega) = \#\{n \geq 1 : X(\omega) \geq n\} =
\sum_{n\geq1}\mathbf{1}_{X \geq n}(\omega)$। द्विक कुल $\bigl(\mathbf{1}_{X \geq n}(\omega)\,\P(\{\omega\})\bigr)_{n,
\omega}$ अऋणात्मक है, अतः कुलों के लिए फ़ूबिनी
(\cref{ch:b2:series}) बिना किसी शर्त के लागू होती है: पहले $n$ में योग लेने पर
$\E(X)$ मिलता है, पहले $\omega$ में योग लेने पर $\sum_n \P(X
\geq n)$; और दोनों एक साथ परिमित
तथा बराबर हैं। $X \sim
\mathcal{G}(p)$ के लिए: $\P(X \geq n) = q^{n-1}$ ($q = 1-p$), अतः $\E(X) =
\sum_{n\geq1}q^{n-1} = \frac{1}{1 - q} = \frac1p$।
\end{solution}

\begin{solution}{exo:b2:randomvar:5}
\emph{सममिति:} खींची गई $i$-वीं गेंद कलश की कोई एकसमान यादृच्छिक
गेंद है ($N$ गेंदों में से किसी के भी खींच-क्रम में स्थान $i$ पर
आने की संभावना बराबर है), अतः $\P(Y_i = 1) = \frac MN =
p$ और रैखिकता से $\E(X) = np$ --- किसी
स्वतंत्रता की ज़रूरत नहीं।

\emph{\omterm{def:b2:randomvar:variance}{सहप्रसरण}:} $i \neq j$ के लिए $\E(Y_iY_j) = \P(\text{खींचें }
i, j \text{ दोनों
सफ़ेद}) = \frac{M(M-1)}{N(N-1)}$ (भिन्न स्थानों के क्रमित युग्मों को भिन्न गेंदों का
क्रमित युग्म \omterm{def:b2:funcseq:def}{एकसमान रूप से} मिलता है)। इसलिए
\[
\operatorname{Cov}(Y_i, Y_j)
= \frac{M(M-1)}{N(N-1)} - \frac{M^2}{N^2}
= \frac{M(N - M)}{N^2}\cdot\frac{-1}{N-1}
= -\frac{p(1-p)}{N-1} < 0 :
\]
अर्थात् कोई सफ़ेद गेंद खींचना बाक़ी खींचों के लिए सफ़ेदों को और विरल
कर देता है। \cref{thm:b2:randomvar:variancerules} से,
\[
V(X) = np(1-p) + n(n-1)\Bigl(-\frac{p(1-p)}{N-1}\Bigr)
= np(1-p)\,\frac{N - n}{N - 1} \leq np(1-p) :
\]
यानी बिना प्रतिस्थापन प्रतिचयन का माध्य वही है पर \omterm{def:b2:randomvar:variance}{प्रसरण}
प्रतिस्थापन वाली स्थिति से \emph{छोटा} है (समता केवल $n = 1$ के लिए),
और ऋणात्मक सहसंबंध किसी स्थिरक की तरह काम करते हैं। $n = N$ के लिए
\omterm{def:b2:randomvar:variance}{प्रसरण} लुप्त हो जाता है: तब गिनती निर्धारणात्मक है।
\end{solution}

\begin{solution}{exo:b2:randomvar:6}
$m = \E(X)$ के परितः प्रसार करने पर:
\[
\E\bigl((X - c)^2\bigr)
= \E\bigl((X - m)^2\bigr) + 2(m - c)\,\E(X - m) + (m - c)^2
= V(X) + (m - c)^2 ,
\]
जो ठीक $c = m$ पर न्यूनतम है और उसका मान $V(X)$ --- अर्थात् माध्य वर्ग के
अर्थ में \omterm{def:b2:randomvar:expectation}{प्रत्याशा} सर्वोत्तम अचर पूर्वानुमानक है।

यदि $\P(X = m) = 1$ हो, तो $(X - m)^2$ प्रायिकता $1$ के साथ लुप्त होता है, अतः $V(X) = 0$
(परिभाषा वाले कुल के पद किसी शून्य समुच्चय के बाहर शून्य हैं)।
विलोमतः, यदि $V(X) = 0$, तो चेबिशेव (\cref{thm:b2:randomvar:markov}) हर $n$ के लिए $\P\bigl(\abs{X - m} \geq
\frac1n\bigr) \leq n^2\,V(X) = 0$ देती है;
घटनाएँ $\bigl\{\abs{X - m} \geq \frac1n\bigr\}$ बढ़कर $\{X \neq
m\}$ तक जाती हैं, अतः एकदिष्ट संततता (\cref{thm:b2:proba:continuity}) $\P(X \neq m) = 0$
दे देती है।
\end{solution}

\begin{solution}{exo:b2:randomvar:7}
$\E(S_n) = \frac n2$ तथा $V(S_n) = \frac n4$।
\emph{मार्कोव:} $\P\bigl(S_n \geq \frac{3n}4\bigr) \leq
\frac{n/2}{3n/4} = \frac23$ --- कोई अचर परिबंध, जो बड़े $n$ के लिए बेकार है।
\emph{चेबिशेव:} \omterm{def:b2:proba:space}{घटना} से $\abs{S_n - \frac n2} \geq
\frac n4$ निकलता है, अतः प्रायिकता $\leq \frac{n/4}{(n/4)^2} =
\frac4n$ है ---
जो क्षय तो करती है, पर केवल बहुपदीय रूप से।
\emph{चेर्नोफ़:} स्वतंत्रता से $\E(e^{tS_n}) =
\prod_{i=1}^n\E(e^{tX_i}) = \bigl(\frac{1 + e^t}{2}\bigr)^n$, और $e^{tS_n} \geq e^{3nt/4}$ पर लगाई गई मार्कोव हर
$t >
0$ के लिए देती है
\[
\P\Bigl(S_n \geq \frac{3n}4\Bigr)
\leq \Bigl(\frac{1 + e^t}{2}\Bigr)^n e^{-3nt/4}
= \exp\Bigl(n\bigl(\ln\tfrac{1 + e^t}{2} - \tfrac{3t}4\bigr)\Bigr).
\]
घातांक को न्यूनतम कीजिए: $e^t = 3$ पर $\frac{\dd}{\dd t}\ln\frac{1+e^t}{2} =
\frac{e^t}{1 + e^t} = \frac34$, अर्थात् $t = \ln 3$, जिससे
\[
\P\Bigl(S_n \geq \frac{3n}4\Bigr)
\leq \Bigl(\frac{4}{2}\Bigr)^n 3^{-3n/4}
= \bigl(2 \cdot 3^{-3/4}\bigr)^n \approx (0.877)^n ,
\]
मिलता है, जो चरघातांकी रूप से छोटा है। पदानुक्रम मार्कोव $\to$
चेबिशेव $\to$ चेर्नोफ़ मानक सीढ़ी है: हर पायदान मार्कोव को चर के
किसी तेज़ी से बढ़ते फलन पर लगाता है।
\end{solution}

\begin{solution}{exo:b2:randomvar:8}
$S_n \sim
\mathcal{B}(n, x)$ वाले $f\bigl(\frac{S_n}{n}\bigr)$ पर लगाई गई अंतरण प्रमेय (\cref{thm:b2:randomvar:transfer}) से:
\[
\E\Bigl[f\Bigl(\frac{S_n}{n}\Bigr)\Bigr]
= \sum_{k=0}^n f\Bigl(\frac kn\Bigr)\binom nk x^k(1-x)^{n-k}
= B_nf(x) .
\]
$\delta > 0$ स्थिर कीजिए और \omterm{def:b2:proba:space}{घटना} $D
= \bigl\{\abs{\frac{S_n}{n} - x} \geq \delta\bigr\}$ पर $\abs{f(S_n/n) - f(x)}$ को बाँटिए: $D$ के बाहर अंतर
अधिक से अधिक संततता का मापांक $\omega_f(\delta) = \sup_{\abs{s - t}\leq\delta}\abs{f(s) -
f(t)}$ है; $D$ पर अधिक से अधिक $2\norm f_\infty$।
\omterm{def:b2:randomvar:expectation}{प्रत्याशा} लेकर और $V\bigl(\frac{S_n}{n}\bigr) =
\frac{x(1-x)}{n} \leq \frac{1}{4n}$ के साथ चेबिशेव बरतकर:
\[
\abs{B_nf(x) - f(x)}
\leq \E\,\abs{f(S_n/n) - f(x)}
\leq \omega_f(\delta)
+ 2\norm f_\infty\,\P(D)
\leq \omega_f(\delta) + \frac{2\norm f_\infty}{4n\delta^2} .
\]
$[0, 1]$ पर $f$ की एकसमान संततता $\omega_f(\delta) \to
0$ बना देती है: पहले $\delta$ चुनिए,
फिर $n$, और \omterm{def:b2:funcseq:def}{एकसमान रूप से} $B_nf \to f$ --- यानी \cref{ch:b2:funcseq} की वाइरश्ट्रास
सन्निकटन प्रमेय, जिसकी ``गणना-प्रमेयिका'' अब द्विपद \omterm{def:b2:randomvar:law}{नियम} के लिए
चेबिशेव असमिका के रूप में पहचानी जा सकती है।
\end{solution}

\begin{solution}{exo:b2:randomvar:9}
$S_n^4 = \sum_{i,j,k,l}X_iX_jX_kX_l$ को खोलिए और \omterm{def:b2:randomvar:expectation}{प्रत्याशा} लीजिए। स्वतंत्रता तथा केंद्रण से, जिस भी
पद में कोई सूचकांक ठीक एक बार आता हो वह लुप्त हो जाता है ($\E(X_i) = 0$
गुणनखंड के रूप में बाहर आ जाता है)। बचे हुए पद: $n$ विकर्ण पद $\E(X_i^4)$,
और बराबर सूचकांकों के दो युग्मों को जोड़ने वाले पद $\E(X_i^2X_j^2) =
\E(X_1^2)^2$ ($i \neq j$ के
लिए), जो $3n(n-1)$ बार आते हैं: मानों का अक्रमित युग्म चुनिए ($\binom n2$
तरीक़े), फिर उन्हें चारों ख़ानों में रखने के $\frac{4!}{2!\,2!} = 6$ तरीक़े --- $6\binom n2 = 3n(n-1)$।
इसलिए $\E(X_1^2)^2
\leq \E(X_1^4)$ (जेनसन या कोशी--श्वार्ज़) के साथ,
\[
\E(S_n^4) = n\,\E(X_1^4) + 3n(n-1)\,\E(X_1^2)^2
\leq C n^2,
\qquad C = 4\,\E(X_1^4) .
\]
कोटि 4 पर मार्कोव:
\[
\P\Bigl(\Bigl|\frac{S_n}{n}\Bigr| \geq \varepsilon\Bigr)
= \P\bigl(S_n^4 \geq n^4\varepsilon^4\bigr)
\leq \frac{Cn^2}{n^4\varepsilon^4}
= \frac{C}{n^2\varepsilon^4} ,
\]
जो कोई \omterm{def:b2:series:summable}{योग्य} श्रेणी है। बोरेल--कांतेली~1 (\cref{thm:b2:proba:borelcantelli}) से हर $j$ के लिए
\omterm{def:b2:proba:space}{घटना} $B_j = \limsup_n\bigl\{\abs{S_n/n} \geq \frac1j\bigr\}$ की प्रायिकता $0$ है, अतः \omterm{def:b2:structures:countable}{गणनीय} उप-योज्यता से $\P\bigl(\bigcup_j B_j\bigr) = 0$।
पूरक पर --- जिसकी प्रायिकता $1$ है --- हर $j$ के लिए ऐसा $N$ है
कि सब $n \geq N$ के लिए $\abs{S_n/n} < \frac1j$: अर्थात् ठीक $S_n/n \to 0$। बृहत् संख्याओं का प्रबल
नियम चतुर्थ आघूर्ण के अंतर्गत टिकता है; उस परिकल्पना को हटाना
(कोल्मोगोरोव की प्रमेय) वर्ष 3 का काम है।
\end{solution}

\begin{solution}{exo:b2:randomvar:10}
$\P(M \leq k) = \bigl(\frac k6\bigr)^2$ (दोनों पासे \omterm{def:b2:proba:independence}{स्वतंत्र} रूप से अधिक से अधिक $k$), अतः $\P(M \geq k) = 1 -
\bigl(\frac{k-1}6\bigr)^2$ और
\[
\E(M) = \sum_{k=1}^6\P(M \geq k)
= 6 - \frac{0 + 1 + 4 + 9 + 16 + 25}{36}
= 6 - \frac{55}{36} = \frac{161}{36} \approx 4.47 ,
\]
जो किसी अकेले पासे के माध्य $3.5$ से आराम से ऊपर है, जैसा किसी
अधिकतम को होना ही चाहिए।
\end{solution}

\begin{solution}{exo:b2:randomvar:11}
$I_i = \mathbf 1_{\sigma(i) = i}$ के साथ: $\P(\sigma(i) = i) =
\frac{(n-1)!}{n!} = \frac1n$, अतः $\E(F_n) = n\cdot\frac1n =
1$। $i \neq j$ के लिए: $\P(\sigma(i) = i, \sigma(j) = j) =
\frac{(n-2)!}{n!} = \frac1{n(n-1)}$, इसलिए
\[
\operatorname{Cov}(I_i, I_j) = \frac1{n(n-1)} - \frac1{n^2}
= \frac{1}{n^2(n-1)} .
\]
\omterm{def:b2:randomvar:variance}{प्रसरण} की औज़ार-पेटी (\cref{thm:b2:randomvar:variancerules}) से,
\[
V(F_n) = n\cdot\frac1n\Bigl(1 - \frac1n\Bigr)
+ n(n-1)\cdot\frac1{n^2(n-1)}
= 1 - \frac1n + \frac1n = 1 .
\]
माध्य $1$, \omterm{def:b2:randomvar:variance}{प्रसरण} $1$, जो $n$ से \omterm{def:b2:proba:independence}{स्वतंत्र} है --- और यह मिलान
समस्या की प्वासों सीमा (\cref{exo:b2:proba:5}) से संगत है।
\end{solution}

\begin{solution}{exo:b2:randomvar:12}
$T_n = \sum_{k=1}^nG_k$, जहाँ $G_k \sim \mathcal G(k/n)$ $k$ के शेष रहने पर किसी नए खिलौने को देखने का समय
है, और चरण \omterm{def:b2:proba:independence}{स्वतंत्र} हैं। इसलिए \cref{ex:b2:fourier:basel} से
\[
V(T_n) = \sum_{k=1}^n\frac{1 - k/n}{(k/n)^2}
\leq \sum_{k=1}^n\frac{n^2}{k^2}
\leq \frac{\pi^2}6\,n^2 ,
\]
$\E(T_n) = nH_n$ के साथ $H_n =
\sum_1^n\frac1k$ (\cref{exo:b2:randomvar:3}), और चेबिशेव $\varepsilon > 0$ के लिए देती है
\[
\P\bigl(\abs{T_n - nH_n} \geq \varepsilon\,n\ln n\bigr)
\leq \frac{\pi^2n^2/6}{\varepsilon^2n^2\ln^2 n}
= \frac{\pi^2}{6\,\varepsilon^2\ln^2n}
\xrightarrow[n\to\infty]{} 0 .
\]
चूँकि $H_n \sim \ln n$, $n\ln n$ से भाग देने पर दिखता है कि प्रायिकता में $T_n/(n\ln n) \to 1$:
अर्थात् $T_n$ के उतार-चढ़ाव कोटि $n$ के हैं, जो माध्य $n\ln n$ के सामने
नगण्य हैं।
\end{solution}

\begin{solution}{pb:b2:randomvar:1}
\textbf{1.} $\E(S_n) = \frac n2$ तथा मार्कोव (\cref{thm:b2:randomvar:markov}) $\P(S_n \geq an) \leq
\frac{n/2}{an} = \frac1{2a}$ देते हैं। मार्कोव केवल
माध्य जानती है: वह $n/2$ पर संकेंद्रित किसी चर को $0$ तथा $n$ के बीच
फैले किसी चर से अलग नहीं कर सकती, अतः वह पुच्छ का मूल्य ऐसे लगाती है
मानो सारा द्रव्यमान वहीं बैठ सकता हो।

\textbf{2.} निष्पक्ष द्विपद $n/2$ के परितः \omterm{def:b2:quadratic:adjoint}{सममित} है ($S_n$ तथा $n - S_n$ का
\omterm{def:b2:randomvar:law}{नियम} एक ही है), अतः $x = n(a -
\frac12) > 0$ के साथ दोनों घटनाएँ $\{S_n - \frac n2 \geq x\}$ तथा $\{S_n - \frac n2 \leq -x\}$ असंयुक्त
और समसंभावी हैं: $\P(S_n \geq an) = \frac12\P(\abs{S_n - \frac n2} \geq x)$। $V(S_n) = \frac n4$ के साथ चेबिशेव:
\[
\P(S_n \geq an)
\leq \frac12\cdot\frac{n/4}{n^2(a - 1/2)^2}
= \frac1{8n(a - 1/2)^2},
\]
जो $a = \frac34$ पर $\frac2n$ है।

\textbf{3.} स्वतंत्रता तथा गुणनफल प्रमेय से $\E(\eu^{tS_n}) = \bigl(\E \eu^{tX_1}\bigr)^n = \bigl(\frac{1
+ \eu^t}2\bigr)^n$। $\eu^{tS_n}$ पर लगाई गई
मार्कोव:
\[
\P(S_n \geq an) \leq \eu^{-tan}\Bigl(\frac{1 +
\eu^t}2\Bigr)^{\!n} = \exp\Bigl(n\bigl(\ln\tfrac{1 +
\eu^t}2 - ta\bigr)\Bigr).
\]
$t$ में घातांक का अवकलज $\frac{\eu^t}{1 + \eu^t}
- a$ है, जो $\eu^t = \frac a{1-a}$ पर लुप्त होता है, अर्थात्
$t^* =
\ln\frac a{1-a} > 0$; वहाँ $\frac{1 + \eu^{t^*}}2 =
\frac1{2(1-a)}$ और घातांक
\[
n\Bigl(-\ln 2 - \ln(1-a) - a\ln\frac a{1-a}\Bigr)
= -n\bigl(\ln2 + a\ln a + (1-a)\ln(1-a)\bigr) = -n\,I(a),
\]
के बराबर है, जहाँ $\intoo{\frac12}1$ पर $I(\frac12) = 0$ तथा $I'(a) = \ln\frac a{1-a} > 0$: $I(a) > 0$। $a = \frac34$ पर: $\eu^{-I(3/4)} = \frac12(\tfrac34)^{-3/4}(\tfrac14)^{-1/4} =
2\cdot3^{-3/4}$,
यानी \cref{exo:b2:randomvar:7} का परिबंध।

\textbf{4.} $n + 1$ संख्याएँ $\binom nja^j(1-a)^{n-j}$ जुड़कर $1$ बनाती हैं, और सबसे बड़ी वह
है जो $j = k = an$ पर है (यहाँ $\mathcal B(n, a)$ का बहुलक $\floor{(n+1)a} = k$ है)। $1$ तक जुड़ने वाली $n + 1$
संख्याओं का अधिकतम कम से कम $\frac1{n+1}$ होता है:
\[
\binom nk a^k(1-a)^{n-k} \geq \frac1{n+1}
\quad\Longrightarrow\quad
\binom nk \geq \frac{a^{-an}(1-a)^{-n(1-a)}}{n+1}
= \frac{\eu^{nH(a)}}{n+1}.
\]
इसलिए $\P(S_n \geq an) \geq \binom{n}{an}2^{-n} \geq
\eu^{n(H(a) - \ln2)}/(n+1) = \eu^{-nI(a)}/(n+1)$: अर्थात् बहुपदीय गुणक $n + 1$ तक चेर्नोफ़ का घातांक ही
सच्चाई है।

\textbf{5.} $n = 100$, $a = \frac34$: मार्कोव $\frac23$; चेबिशेव $\frac2{100} = 0.02$; चेर्नोफ़ $(2\cdot3^{-3/4})^{100} = \eu^{-100\,I(3/4)} \approx
2.1\cdot10^{-6}$,
जबकि सटीक मान $2.8\cdot10^{-7}$ है। सार: सूचना का हर आघूर्ण परिबंध को बहुपदीय
रूप से बाँट देता है; चरघातांकी आघूर्ण उसकी \emph{प्रकृति} ही बदल
देता है।

\textbf{6.} $\cosh t = \sum_{k\geq0}\frac{t^{2k}}{(2k)!}$ तथा $\eu^{t^2/2} = \sum_{k\geq0}\frac{t^{2k}}{2^kk!}$; दावा पद-दर-पद $(2k)! \geq 2^kk!$ से निकलता है, जो आगमन
से टिकता है: $(2k)! = 2k(2k-1)\cdot(2k-2)! \geq 2k\cdot
2^{k-1}(k-1)! = 2^kk!\cdot(2k-1) \geq 2^kk!$।

\textbf{7.} स्वतंत्रता से $\E\bigl(\eu^{t\sum\varepsilon_i}
\bigr) = (\cosh t)^n \leq \eu^{nt^2/2}$, अतः मार्कोव $\P(\sum\varepsilon_i \geq s) \leq \eu^{nt^2/2 - ts}$ देती है; और $t = s/n$ पर
न्यूनतम करने से $\eu^{-s^2/(2n)}$ मिलता है।

\textbf{8.} $X_i = \frac{1 + \varepsilon_i}2$ के साथ $\widehat
p_n - \frac12 = \frac1{2n}\sum\varepsilon_i$, अतः $\{\widehat p_n - \frac12 \geq \delta\} =
\{\sum\varepsilon_i \geq 2n\delta\}$, और प्रश्न 7 परिबंध $\eu^{-(2n\delta)^2/(2n)} = \eu^{-2n\delta^2}$ दे
देता है। सममित घटना का भी वही परिबंध है, जिससे $\abs{\widehat p_n - \frac12} \geq \delta$ के लिए गुणक $2$
आता है।

\textbf{9.} $\E(\eu^{tX}) = \sum_x\eu^{tx}\P(X = x)$ $t$ के चिकने फलनों की कोई श्रेणी है, जिसके
पद-दर-पद अवकलज हर \omterm{def:b2:metric:compact}{संहत} $t$-अंतराल पर $\eu^{\abs t}\P(X = x)$ से प्रभावित हैं ($0 \leq x \leq 1$):
अतः प्रसामान्य अभिसारी श्रेणियों के लिए अवकलन प्रमेय (\cref{thm:b2:funcseq:differentiation}) से वह
दो बार \omterm{def:b2:diffcalc:differential}{अवकलनीय} है, और विभाग नियम $\psi' = \E_t(X)$ तथा $\psi'' = \E_t(X^2) - \E_t(X)^2$ देता है, जहाँ $\E_t$
पुनःभारित भारों $\eu^{tx}\P(X{=}x)/\E(\eu^{tX})$ के लिए \omterm{def:b2:randomvar:expectation}{प्रत्याशा} है --- जो अऋणात्मक हैं, $1$
तक जुड़ते हैं, और उन्हीं मानों $x \in \intcc01$ द्वारा ढोए जाते हैं। $\intcc01$-मान
वाले किसी चर का \omterm{def:b2:randomvar:variance}{प्रसरण} अधिक से अधिक $\frac14$ होता है: \cref{exo:b2:randomvar:6} से वह $\min_c\E_t((X - c)^2) \leq
\E_t\bigl((X - \tfrac12)^2\bigr) \leq \tfrac14$
है। $\psi(0) = 0$, $\psi'(0) = p$ का उपयोग करते हुए समाकल शेषफल वाली टेलर:
\[
\psi(t) = tp + \int_0^t(t - s)\,\psi''(s)\,\dd s
\leq tp + \frac{t^2}2\cdot\frac14,
\]
अर्थात् सब वास्तविक $t$ के लिए $\E(\eu^{t(X - p)}) \leq \eu^{t^2/8}$।

\textbf{10.} स्वतंत्रता से $\E\bigl(\eu^{t(S_n -
np)}\bigr) \leq \eu^{nt^2/8}$; मार्कोव तथा इष्टतमीकरण $t
= 4\delta$ देते
हैं
\[
\P(\widehat p_n - p \geq \delta)
\leq \eu^{nt^2/8 - tn\delta}\Big|_{t = 4\delta}
= \eu^{-2n\delta^2};
\]
और इसे चरों $1 - X_i$ पर लगाने से (जो भी $\intcc01$ में हैं) दूसरी पुच्छ परिबद्ध
हो जाती है, जिससे द्विपक्षीय $2\eu^{-2n\delta^2}$ मिलता है।

\textbf{11.} चेबिशेव को केवल द्वितीय आघूर्ण चाहिए और वह $\frac{p(1-p)}{n\delta^2}$ देती
है; हॉफडिंग को \emph{परिबद्धता} चाहिए और वह $2\eu^{-2n\delta^2}$ देती है। $p =
\frac12$, $\delta = 0.03$
पर परिबंध (लगभग) $\frac{278}{n}$ बनाम $2\eu^{-0.0018n}$ हैं; वे $n \approx 1200$ के पास कटते हैं, जिसके
बाद चरघातांकी परिबंध जीतता है, और ख़ूब जीतता है ($n = 5000$: $0.056$ बनाम
$2.5\cdot10^{-4}$)।

\textbf{12.} हॉफडिंग (प्रश्न 10) से $2n\delta^2 \geq \ln\frac2\alpha$ होते ही $\P(\abs{\widehat
p_n - p} \geq \delta) \leq 2\eu^{-2n\delta^2} \leq \alpha$, अर्थात् $n \geq
\frac{\ln(2/\alpha)}{2\delta^2}$।

\textbf{13.} $\delta = 0.03$, $\alpha = 0.05$: $n \geq
\frac{\ln 40}{2\cdot0.0009} \approx 2049.4$: $2050$ लोग। $\delta = 0.01$ के लिए: $n \geq \frac{\ln40}{0.0002} \approx
18\,445$। जनसंख्या का
आकार कभी नहीं आता, क्योंकि प्रतिचयित हर मतदाता को कोई नई
बर्नूली$(p)$ खींच माना जाता है: मतसर्वेक्षण की कठिनाई किसी सिक्के का
\omterm{def:b2:randomvar:variance}{प्रसरण} है, देश का आकार नहीं। सीमांत आधा करने पर प्रतिदर्श चार गुना
पड़ता है --- यही $1/\delta^2$ \omterm{def:b2:randomvar:law}{नियम} है।

\textbf{14.} चेबिशेव: $n \geq
\frac1{4\alpha\delta^2}$ के लिए $\P(\abs{\widehat p_n - p} \geq
\delta) \leq \frac{p(1-p)}{n\delta^2} \leq
\frac1{4n\delta^2} \leq \alpha$, अर्थात् तीन अंकों पर $5556$ ---
जो हॉफडिंग की माँग का लगभग $2.7$ गुना है। बिना प्रतिस्थापन \omterm{def:b2:randomvar:variance}{प्रसरण}
$\frac{N -
n}{N-1} < 1$ से गुणित हो जाता है (\cref{exo:b2:randomvar:5}), अतः वही $n$ केवल बेहतर ही कर सकता
है: प्रतिस्थापन वाला परिकलन रूढ़िवादी परिकलन है।

\textbf{15.} $\{\widehat p_n \leq \frac12\} \subseteq
\{\widehat p_n - 0.52 \leq -0.02\}$, अतः एकपक्षीय हॉफडिंग परिबंध से $n \geq \frac{\ln
100}{2\cdot0.0004} \approx 5756.5$ होते ही $\P(\widehat p_n \leq \tfrac12) \leq
\eu^{-2n(0.02)^2} \leq 0.01$:
$5757$ मतदाता। क़ीमत वांछित परिशुद्धता के नहीं, बल्कि \emph{बढ़त} के
व्युत्क्रम वर्ग की तरह मापित होती है: काँटे की टक्कर महँगी पड़ती है।

\textbf{16.} बरता गया: प्रतिदर्श मतदाता-समूह से एकसमान तथा
\omterm{def:b2:proba:independence}{स्वतंत्र} रूप से खींचा जाता है; प्रतिचयित हर व्यक्ति ईमानदारी से
उत्तर देता है, और सर्वेक्षण के दौरान $p$ हिलता नहीं। असली
सर्वेक्षण तीनों का उल्लंघन करते हैं: पहुँच \omterm{def:b2:series:summable}{योग्य} तथा इच्छुक
उत्तरदाता कोई एकसमान प्रतिदर्श नहीं होते (चयन तथा अनुत्तर अभिनति),
और उत्तर असत्य या अस्थिर हो सकते हैं। ये \emph{अभिनति} त्रुटियाँ
हैं: वे $\E(\widehat p_n)$ को $p$ से इतना खिसका देती हैं जो $n$ से \omterm{def:b2:proba:independence}{स्वतंत्र} है,
अतः कोई भी प्रतिदर्श-आकार उन्हें घटाता नहीं --- इस भाग का गणित केवल
उतार-चढ़ाव वाले पद को साधता है।

\textbf{17.} आकार $m$ के किसी एक समूह के लिए चेबिशेव: $m \geq \frac2{\delta^2}$ के लिए
$\P(\abs{
\widehat p^{(i)} - p} \geq \delta) \leq \frac{1}{4m\delta^2}
\leq \frac18$। यदि $k/2$ से कम समूह चूकें, तो मानों $\widehat p^{(i)}$ में से $k/2$ से अधिक
\omterm{def:b2:metric:topology}{विवृत} अंतराल $\intoo{p -
\delta}{p + \delta}$ में पड़ते हैं, और उनकी माध्यिका भी; इसलिए $\{\abs{M - p} \geq \delta\}$ $k$
\omterm{def:b2:proba:independence}{स्वतंत्र} समूहों में कम से कम $\lceil
k/2\rceil$ चूकों को बाध्य कर देता है। चूकने
वाले समूहों के $\binom k{\lceil k/2\rceil}$ संभव समुच्चयों पर संघ परिबंध देता है
\[
\P(\abs{M - p} \geq \delta)
\leq \binom{k}{\lceil k/2\rceil}
\Bigl(\frac18\Bigr)^{k/2}
\leq 2^k\,8^{-k/2} = 2^{-k/2} :
\]
अर्थात् समूहों की संख्या में चरघातांकी क्षय, जो प्रसरणों के सिवा
कुछ भी बरते बिना ख़रीदा गया --- और यह ठीक तब उपयोगी है जब पद
अपरिबद्ध हों और हॉफडिंग उपलब्ध न हो।

\textbf{18.} कोशी--श्वार्ज़ (\cref{thm:b2:randomvar:jensen}):
\[
\E(X) = \E(X\,\mathbf 1_{X>0})
\leq \sqrt{\E(X^2)}\sqrt{\E(\mathbf 1_{X>0}^2)}
= \sqrt{\E(X^2)\,\P(X > 0)} ;
\]
वर्ग कीजिए और भाग दीजिए।

\textbf{19.} मान लीजिए $h(a) = I(a) - 2(a - \tfrac12)^2$। तब $h(\tfrac12) = 0$, $h'(a) = \ln\frac a{1-a} - 4(a -
\tfrac12)$ $\tfrac12$ पर लुप्त होता है, और
\[
h''(a) = \frac1a + \frac1{1-a} - 4 = \frac{1}{a(1-a)} - 4
\geq 0
\]
क्योंकि $a(1-a) \leq \frac14$। अतः $h'$ $\intco{\frac12}1$ पर $0$ से बढ़ता है, इसलिए $h' \geq 0$ तथा $h \geq 0$:
$I(a)
\geq 2(a - \tfrac12)^2$।

\textbf{20.} $I(\tfrac12) = I'(\tfrac12) = 0$, $I''(a) =
\frac1{a(1-a)}$ $I''(\tfrac12) = 4$ देता है, और $I'''(\tfrac12)
= 0$ (फलन $\tfrac12$ के परितः \omterm{def:b2:quadratic:adjoint}{सममित}
है), अतः $I(\tfrac12 + \delta) = 2\delta^2 + O(\delta^4)$। तब प्रश्न 4 सच्ची पुच्छ को \emph{नीचे} से $\eu^{-n(2\delta^2 + O(\delta^4))}/(n+1)$ द्वारा
परिबद्ध कर देता है: अर्थात् छोटे $\delta$ के लिए हॉफडिंग घातांक $2n\delta^2$
अनंतस्पर्शी रूप से सटीक है --- केवल $n$ में बहुपदीय सुधार संभव हैं।

\textbf{21.} मार्कोव: एक आघूर्ण, क्षय $1/a$, जो केवल बाक़ी सबके पीछे
के इंजन के रूप में उपयोगी है (प्रश्न 1 उसे सपाट दिखाता है)।
चेबिशेव: दो आघूर्ण, क्षय $\frac{V}{n\delta^2}$, जो बिना किसी अतिरिक्त परिकल्पना के
तीखी है (\cref{ex:b2:randomvar:chebsharp}), और प्रश्न 23 पर सर्वोत्तम औज़ार। चतुर्थ आघूर्ण
(\cref{exo:b2:randomvar:9}): क्षय $C/n^2$, जो किसी प्रबल नियम में सिकुड़ने के लिए बस उतनी
ही योग्यता देता है। हॉफडिंग: परिबद्ध चर, क्षय $2\eu^{-2n\delta^2}$, भाग III का
कर्मठ घोड़ा। सटीक दर $I(a)$ वाली चेर्नोफ़: पूरे चरघातांकी आघूर्ण,
अजेय घातांक (प्रश्न 4, 20), और बाक़ी सबका संदर्भ बिंदु।

\textbf{22.} यदि सिक्का निष्पक्ष हो: $\P(\widehat p_n > 0.525)
\leq \P(\widehat p_n - \tfrac12 \geq 0.025) \leq
\eu^{-2n(0.025)^2}$। यदि $p = 0.55$: $\P(\widehat p_n \leq
0.525) \leq \P(\widehat p_n - 0.55 \leq -0.025) \leq
\eu^{-2n(0.025)^2}$। दोनों
त्रुटियाँ $2n(0.025)^2 \geq \ln 100$ होने पर $0.01$ से नीचे हैं, अर्थात् $n \geq 3684.2$: $3685$ उछाल।
($2.5$ अंक दूर की परिकल्पनाओं को अलग करने की क़ीमत उतनी ही है जितनी
$\pm2.5$ अंकों तक आकलन करने की।)

\textbf{23.} हॉफडिंग: $n \geq \frac{\ln 40}{2(0.005)^2}
\approx 73\,778$। सच्चे \omterm{def:b2:randomvar:variance}{प्रसरण} $p(1-p)
= 0.0099$ वाली चेबिशेव: $n \geq \frac{0.0099}{0.05\cdot(0.005)^2} =
7920$ ---
नौ गुना सस्ती। हॉफडिंग का घातांक $2n\delta^2$ \omterm{def:b2:randomvar:variance}{प्रसरण} का मूल्य उसकी सबसे बुरी
स्थिति $\frac14$ पर लगाता है, जो $p = 0.01$ होने पर बेतुका निराशावादी है; विनम्र
द्वितीय आघूर्ण बेहतर जानता है। छूटा हुआ औज़ार कोई प्रसरण-सजग
चरघातांकी परिबंध है (बर्नस्टाइन असमिका, वर्ष 3) --- अथवा, विरल
\omterm{def:b2:proba:space}{घटनाओं} के लिए, \cref{ch:b2:genfun} में सिद्ध प्वासों सन्निकटन, जो स्वाभाविक सापेक्ष
मापक्रम पर काम करता है।

\textbf{24.} $\delta > 0$ स्थिर कीजिए: $\sum_n 2\eu^{-2n\delta^2} <
\infty$ (ज्यामितीय-प्रकार की श्रेणी), अतः
बोरेल--कांतेली~1 (\cref{thm:b2:proba:borelcantelli}) देती है $\P(\abs{\widehat p_n - p} \geq \delta \text{ अपरिमित
बार}) = 0$, अर्थात् हर $j$ के लिए \omterm{def:b2:proba:space}{घटना} $E_j =
\bigcup_N\bigcap_{n\geq N}\{\abs{\widehat p_n - p} <
\tfrac1j\}$ की प्रायिकता $1$ है।
\omterm{def:b2:structures:countable}{गणनीय} प्रतिच्छेद $\bigcap_jE_j$ की प्रायिकता भी $1$ है (पूरकों पर
उप-योज्यता), और उस पर $\widehat p_n
\to p$: अर्थात् सिक्का-उछालों के लिए बृहत्
संख्याओं का प्रबल नियम, जहाँ परिबद्धता वही भूमिका निभा रही है जो
\cref{exo:b2:randomvar:9} में चतुर्थ आघूर्ण ने निभाई थी।

\textbf{25.} एक आघूर्ण कोई सपाट परिबंध ख़रीदता है; दो $1/(n\delta^2)$ ख़रीदते
हैं, और उससे अधिक नहीं (तीखेपन का उदाहरण); चार $1/n^2$ ख़रीदते हैं, जो
किसी लगभग-निश्चित \omterm{def:b2:randomvar:law}{नियम} में सिकुड़ने के लिए पर्याप्त है; परिबद्धता
$\eu^{-2n\delta^2}$ ख़रीदती है; और पूरा चरघातांकी आघूर्ण सटीक दर $I$ ख़रीदता है,
जिसे कोई विधि नहीं हरा सकती। $2050$ लोगों का सर्वेक्षण किसी भी देश के
लिए इसलिए पर्याप्त है कि प्रतिदर्श का उतार-चढ़ाव सिक्के के \omterm{def:b2:randomvar:variance}{प्रसरण}
से चलता है, जनसंख्या के आकार से नहीं --- $1/\delta^2$ तथा $\ln(1/\alpha)$ वाली क़ीमत की
पर्चियाँ सार्वभौमिक हैं। वर्ष 3 के खंड की केंद्रीय सीमा प्रमेय इन
असमिकाओं को $\sqrt n$ मापक्रम पर स्पष्ट अचरों वाले किसी सटीक सीमा नियम
से बदल देती है --- और इस तरह इस समस्या के हर परिबंध को किसी
अनंतस्पर्शी समता में बदल देती है।
\end{solution}
